\documentclass[11pt]{article}
\usepackage[margin=1in]{geometry}
\usepackage{amsmath,amssymb,amsthm,booktabs,hyperref}
\newtheorem{theorem}{Theorem}
\newtheorem{lemma}{Lemma}
\title{No Higgs $\mathbf{10}$ on the $e_6+a_2$ route of the non-supersymmetric heterotic string\\[2pt]
\large A lattice theorem for $T^6/\mathbb Z_3$ orbifolds of the $SO(16)\times SO(16)$ string}
\author{W(3,3) programme --- Claude track (Passes 12161--12252)}
\date{October 2026}
\begin{document}
\maketitle

\begin{abstract}
In $T^6/\mathbb Z_3$ orbifolds of the $SO(16)\times SO(16)$ heterotic string built on a three-family $e_6+a_2$ shift, Wilson-line breaking
to $SU(5)$ leaves fractional charges (Wen--Witten), and the field-theoretic escape needs a scalar in the $\mathbf{10}$ of $SU(5)$. We show
that no such scalar exists for any Wilson-line configuration that leaves an $SU(5)$, in any sector. The proof has two parts: a statistics
rule, derived from the projection phase of the fermion-parity element, that decides boson versus fermion from the left-moving momentum
alone; and an exhaustive enumeration, by coset tables over $E_8/3E_8$, of the twisted massless $\mathbf{10}$-type momenta. Twisted
$\mathbf{10}$s do occur, in $0.18\%$ of one-line and $1.8\%$ of two-line configurations, but every one has the fermionic value of the rule.
\end{abstract}

\section{Setup}
The orbifold group is $\mathbb Z_2^W\times\mathbb Z_3$. The fermion-parity (Witten) element $g=(-1)^F e^{2\pi i\,p\cdot V_0}$ acts
trivially on $T^6$; its gauge shift $V_0$ is a \emph{split} involution of $E_8\times E_8$, $V_0^2=2$, the class that makes the
ten-dimensional parent the tachyon-free $SO(16)\times SO(16)$ string. The $\mathbb Z_3$ shift $V_1$ is one of seven $e_6+a_2$ bases with
three chiral families; torus one carries a family Wilson line $W_f$; tori two and three carry order-three Wilson lines $W_1,W_2$.
Sectors are labelled by the Witten label $w\in\{0,1\}$, the $\mathbb Z_3$ power $n$ and fixed-point labels $m_a$.

\section{The statistics rule}
\begin{lemma}\label{lem:rule}
On a state of sector $h=(w,n,m_a)$ with left momentum $P$ and right-moving weight $R$, invariance under $g$ requires
\[
P\cdot V_0-R_0\;\equiv\;\tfrac12\Bigl(w\,(V_0^2-1)+n\,V_1\cdot V_0+\textstyle\sum_a m_a\,W_a\cdot V_0\Bigr)\pmod 1 ,
\]
$R_0$ the spacetime component of $R$ (integral for bosons, half-integral for fermions). If $V_0^2=2$ and all products of $V_0$ with
the shifts and Wilson lines are even, the right side is $w/2$, and a state is a boson iff $P\cdot V_0\equiv w/2$.
\end{lemma}
\begin{proof}
The phase of $g$ on the state is $(P\cdot V_0-R\cdot s_0)-\tfrac12(V_h\cdot V_0-v_h\cdot s_0)$; $g$ is a pure shift, so there are no
oscillator or fixed-point phases. With $R\cdot s_0=R_0$, $v_h\cdot s_0=w$ and $V_h=wV_0+nV_1+\sum m_aW_a$ the formula follows.
\end{proof}
\emph{Consistency forces the evenness.} If some $W\cdot V_0$ is odd, the right side depends on whether a fixed-point label is written
$1$ or $-2$; the projection is ill-defined. The orbifolder rejects every such configuration ($147/147$), and $65{,}659$ of the $231{,}124$
one-line configurations that satisfy only $W\cdot V_0\in\mathbb Z$ are of this kind.
\emph{Check.} From left weights alone the rule reproduces the engine's statistics for all $5{,}974$ fields of eight full $T^6/\mathbb Z_3$
spectra and for all $74{,}174$ fields of the $104$ $A_8$ Standard Models of the companion route.

\section{The theorem}
\begin{theorem}[no Higgs $\mathbf{10}$]
For every consistent configuration of order-three Wilson lines that leaves an $SU(5)$, no state in the $\mathbf{10}$ of $SU(5)$ is a scalar.
\end{theorem}
\begin{proof}[Proof (computer-assisted)]
\emph{Untwisted:} the $\mathbf{10}$ lies in the $\mathbf{16}$ of $SO(10)$, with $P\cdot V_0=\tfrac12$, $w=0$: fermions (all $7\times864$ lines).
\emph{Twisted:} a massless twisted state has $|P|^2=4/3$ (with an oscillator the shell $2/3$ is below the $\mathbf{10}$-weight norm $6/5$);
a $\mathbf{10}$ needs the $SU(5)$ projection of the first-$E_8$ part to have norm $6/5$. Commuting order-three lines diagonalise on the
$D_5$ Cartan, so an $SU(5)$ intersection forces one line of $SU(5)$ type and the other preserving it. For each base, the cosets
$V_1+n_1W_f+wV_0+n_2\lambda_1+n_3\lambda_2$ (and their second-$E_8$ partners) are fixed by classes in $E_8/3E_8$, so all short vectors are
tabulated once and every pair is decided by table look-up. Results:
\begin{center}\begin{tabular}{lrrr}
\toprule
& consistent configurations & with a twisted $\mathbf{10}$ & scalar type\\
\midrule
one line & $165{,}465$ & $294$ & $0$\\
two lines & $119{,}915{,}126$ & $2{,}210{,}046$ & $0$\\
\bottomrule
\end{tabular}\end{center}
Every twisted $\mathbf{10}$-type momentum of the one-line class lies in a Witten-twisted coset with $P\cdot V_0\in\mathbb Z$; by the lemma
all are fermions. Positive controls: an independent random sample of $366$ two-line hits is recovered exactly; bases with no hits in the
sample have none in the enumeration.
\end{proof}

\section{Caveats}
\begin{itemize}
\item \emph{Wilson-line consistency.} The correct conditions are those of Pl\"oger, Ramos-S\'anchez, Ratz and Vaudrevange
(hep-th/0702176, eqs.~(7a)--(7f)), derived from $\Phi_{\rm vac}(g,h)^{\gcd(n,s)}=1$ (their (A.11)): in particular $\gcd(N_a,N_b)\,(W_a\cdot W_b)
\in2\mathbb Z$ and $N_a(W_a\cdot V_i)\in2\mathbb Z$ --- the latter, for $V_0$, is our parity condition. The non-supersymmetric orbifolder does not impose the first (its failure flag sits in a verbose-only branch). A patched build flags $497$ corpus inputs; every one has a lattice-shift repair with a spectrum identical to the stock output ($983/983$), and at order three the lattice representative carries no torsion ($160/160$ relabellings invariant). The theorem is a statement of absence and holds in any case.
\item \emph{Engine dependence, resolved.} The two $N=0$ engines agree field for field in the untwisted and $\mathbb Z_3$-twisted sectors; the orbifolder~1.2.1-based build drops every Witten-twisted fixed brane (a local condition it applies there fails), which is where all the extra twisted $\mathbf{10}$s of the non-supersymmetric engine sit.
\item \emph{Scope.} One geometry ($T^6/\mathbb Z_3$), one shift class ($e_6+a_2$), shift orbifolds only.
\end{itemize}

\paragraph{Consequence.} Wilson-line breaking on this route leaves fractional charges, and the field-theoretic escape is closed. Together
with the Weinberg-angle and proton obstructions recorded elsewhere, this route does not give a realistic Standard Model.

\paragraph{Reproducibility.} Producers \texttt{w33\_pass12217\ldots}, \texttt{12229\ldots}, \texttt{12235\ldots}, \texttt{12241\ldots},
\texttt{12247\ldots}, \texttt{12253\ldots} in \texttt{analysis/}, with drivers in \texttt{analysis/search\_drivers\_122*}.
\end{document}
